\documentclass[11pt,a4paper,oneside]{book}
\usepackage{geometry}
\geometry{top=2.4cm,bottom=2.5cm,left=2.2cm,right=2.2cm,headheight=16pt}
\usepackage{amsmath,amssymb,mathtools,bm}
\usepackage{booktabs,longtable,tabularx,array,multirow}
\usepackage{graphicx}
\usepackage{tikz}
\usetikzlibrary{arrows.meta,positioning,shapes.geometric,calc}
\usepackage{xcolor}
\usepackage{fancyhdr}
\usepackage{hyperref}
\hypersetup{colorlinks=true,linkcolor=blue!55!black,urlcolor=blue!55!black,citecolor=blue!55!black,pdftitle={آموزش جامع تحلیل قابلیت اطمینان و جایابی تجهیزات حفاظتی شبکه توزیع},pdfauthor={OpenAI}}
\usepackage{xepersian}
\settextfont[Scale=1.03]{Vazirmatn}
\setlatintextfont{DejaVu Serif}
\setdigitfont[Scale=1.0]{Vazirmatn}
\setmonofont{DejaVu Sans Mono}
\renewcommand{\labelitemi}{\ensuremath{\bullet}}
\linespread{1.24}
\setlength{\parskip}{4pt}
\setlength{\parindent}{1.2em}
\renewcommand{\chaptername}{فصل}
\renewcommand{\contentsname}{فهرست مطالب}
\renewcommand{\listfigurename}{فهرست شکل‌ها}
\renewcommand{\listtablename}{فهرست جدول‌ها}
\pagestyle{fancy}
\fancyhf{}
\fancyhead[R]{\small تحلیل قابلیت اطمینان و جایابی تجهیزات شبکه توزیع}
\fancyhead[L]{\small مطالعه موردی فیدرهای پست کوثر}
\fancyfoot[C]{\thepage}
\setcounter{secnumdepth}{3}
\setcounter{tocdepth}{2}
\newcommand{\file}[1]{\lr{\texttt{#1}}}
\newcommand{\code}[1]{\lr{\texttt{#1}}}
\newcommand{\unit}[1]{\,\mathrm{#1}}
\newcommand{\summary}[1]{\subsection*{جمع‌بندی فصل #1}}
\newcommand{\exercise}{\subsection*{تمرین کوتاه و پاسخ}}
\newcommand{\answer}{\noindent\textbf{پاسخ:} }
\newcommand{\warn}{\noindent\textbf{نکته مهم داده‌ای:} }
\newcommand{\assumption}{\noindent\textbf{فرض آموزشی:} }
\newcommand{\Rrule}[1]{\lr{\texttt{#1}}}

\begin{document}
\begin{titlepage}
\centering
\vspace*{1.3cm}
{\Large راهنمای آموزشی مهندسی برق قدرت و شبکه‌های توزیع\par}
\vspace{1.2cm}
{\Huge\bfseries تحلیل قابلیت اطمینان، محاسبه \lr{ENS} و \lr{EENS} و جایابی بهینه تجهیزات حفاظتی\par}
\vspace{0.8cm}
{\LARGE با فرمول‌بندی کامل ریاضی و مطالعه موردی فیدرهای پست کوثر\par}
\vspace{1.3cm}
\begin{tikzpicture}[node distance=1.3cm,>=Stealth, every node/.style={draw,rounded corners,align=center,minimum height=0.9cm,minimum width=2.5cm,font=\small}]
\node (data) {داده خاموشی\ و \lr{GIS}};
\node[right=of data] (clean) {پاک‌سازی و\ طبقه‌بندی};
\node[right=of clean] (rel) {شاخص‌های\ قابلیت اطمینان};
\node[below=of rel] (eens) {محاسبه\ \lr{EENS}};
\node[left=of eens] (opt) {جایابی\ کلیدها};
\node[left=of opt] (eval) {ارزیابی و\ تحلیل حساسیت};
\draw[->,thick] (data)--(clean);
\draw[->,thick] (clean)--(rel);
\draw[->,thick] (rel)--(eens);
\draw[->,thick] (eens)--(opt);
\draw[->,thick] (opt)--(eval);
\end{tikzpicture}
\vfill
{\large نسخه آموزشی مبتنی بر داده‌های پیوست‌شده\par}
{\large فایل‌های خاموشی \file{3.xlsx} تا \file{7.xlsx} و لایه‌های \lr{GIS} پروژه \file{Kosar.zip}\par}
\vspace{0.7cm}
{\large آبان ۱۴۰۵\par}
\end{titlepage}

\frontmatter
\chapter*{درباره این راهنما}
این راهنما برای دانشجویان، کارشناسان تازه‌کار و مهندسانی تدوین شده است که می‌خواهند از داده خاموشی و اطلاعات مکانی شبکه، به شاخص‌های قابلیت اطمینان و سپس به تصمیم مهندسی درباره محل نصب ریکلوزر، سکشنالایزر و سکسیونر برسند. مسیر آموزش از داده واقعی آغاز می‌شود؛ سپس تبدیل واحد، پاک‌سازی، طبقه‌بندی رخداد، مدل نرخ خرابی، محاسبه انرژی تأمین‌نشده و در نهایت مسئله بهینه‌سازی تشریح می‌شود.

مطالعه موردی اصلی، فایل‌های واقعی خاموشی فیدرهای پست کوثر است که در آرشیو \file{kosar.rar} قرار داشتند. آرشیو \lr{GIS} درون آن، با نام \file{Kosar.zip} و پوشه خروجی \file{Export\_\allowbreak 1404\_\allowbreak 08\_\allowbreak 17\_\allowbreak 12\_\allowbreak 26\_\allowbreak 05} نیز برای استخراج ویژگی‌های فیدرها استفاده شده است. کد و مستندات الگوریتمی نرم‌افزار از بسته زیر بررسی شده‌اند:
\begin{center}
\file{SwitchLoc-arena-fad3d399-v1.22.0.zip}
\end{center}

\warn نتایج عددیِ جدول‌های تحلیل خاموشی، حاصل اجرای مسیر خواندن و طبقه‌بندی پروژه روی نسخه‌های پیوست‌شده فایل‌هاست. در مقابل، هرجا بار اکتیو، ضریب بار، ضریب توان، ظرفیت پشتیبان یا زمان تعمیر در فایل‌های خام موجود نبوده، مقدار آن به‌عنوان «فرض آموزشی» معرفی شده است؛ این فرض‌ها نباید با اندازه‌گیری واقعی اشتباه گرفته شوند.

\chapter*{راهنمای استفاده}
برای یادگیری پیوسته، فصل‌ها را به‌ترتیب بخوانید. فصل ۲ مشخص می‌کند داده‌ها چه چیزی را بیان می‌کنند و چه چیزی را بیان نمی‌کنند. فصل‌های ۳ و ۴ مبنای محاسبه نرخ خاموشی و مدت خاموشی هستند. فصل ۵ به \lr{EENS} می‌رسد و فصل‌های ۶ تا ۸، ساختار شبکه، انواع تجهیزات و منطق جایابی را شرح می‌دهند. فصل ۹ یک مثال عددی یکپارچه از فیدر \code{F5} ارائه می‌کند. در تمام متن، شماره معادلات با دستور \verb|\eqref{...}| و ارجاع به فصل‌ها با شماره بخش آمده است.

\tableofcontents
\listoftables
\mainmatter

\chapter{صورت مسئله و تصویر کلی مهندسی}
\section{چرا جایابی تجهیزات حفاظتی یک مسئله قابلیت اطمینان است؟}
شبکه توزیع برق معمولاً به‌صورت چند فیدر شعاعی یا نزدیک به شعاعی بهره‌برداری می‌شود. وقوع خطا می‌تواند باعث عملکرد بریکر پست، ریکلوزر، فیوز یا سایر تجهیزات حفاظتی شود. در شبکه‌ای که کلیدهای مانور و نقاط جداسازی مناسبی دارد، ممکن است پس از تشخیص محل خطا فقط بخش معیوب جدا شود و قسمت سالم شبکه با مانور مناسب دوباره برقدار گردد. در شبکه‌ای بدون نقاط جداسازی کافی، یک خطای محدود می‌تواند مشترکان بیشتری را برای مدت طولانی‌تری بی‌برق کند.

هدف از جایابی تجهیزات این نیست که صرفاً تعداد کلیدها زیاد شود. هدف آن است که با تعداد و هزینه محدود تجهیزات، انرژی تأمین‌نشده مورد انتظار کاهش یابد، ایمنی و هماهنگی حفاظتی حفظ شود و امکان مانور واقعی شبکه وجود داشته باشد. از این رو، مسئله به سه لایه وابسته است: (۱) مدل رخدادهای گذشته، (۲) مدل توپولوژی و بار شبکه، و (۳) مدل عملکرد تجهیزات و زمان بازیابی.

\section{زنجیره محاسباتی}
\begin{enumerate}
\item داده خاموشی به رکوردهای استاندارد و قابل ممیزی تبدیل می‌شود.
\item رخدادها بر اساس علت، تجهیز عمل‌کرده و شرح حادثه طبقه‌بندی می‌شوند.
\item با توجه به مدت پوشش معتبر داده، نرخ رخداد و مدت متوسط خاموشی محاسبه می‌شود.
\item مدل شبکه به بخش‌ها یا زون‌ها تقسیم می‌شود و بار هر بخش برآورد یا از اطلاعات مشترکین استخراج می‌گردد.
\item برای هر محل خطا، مدت خاموشی بار معیوب و بار سالمِ قابل بازیابی محاسبه می‌شود.
\item محل و نوع تجهیزات با کمینه کردن \lr{EENS} و رعایت محدودیت‌ها انتخاب می‌شود.
\end{enumerate}

\section{تمایز سه مفهوم مهم}
«تعداد خاموشی» تعداد رخدادهای ثبت‌شده است. «نرخ خاموشی» تعداد رخداد به‌ازای واحد زمان مواجهه است. «انرژی تأمین‌نشده» مقدار انرژی است که به دلیل خاموشی به مصرف‌کننده تحویل نشده است. این سه مفهوم قابل جایگزینی با یکدیگر نیستند. برای مثال، فیدری با تعداد رخداد کمتر اما بار بیشتر یا زمان بازیابی طولانی‌تر ممکن است \lr{EENS} بزرگ‌تری داشته باشد.

\summary{۱}
تحلیل جایابی، زنجیره‌ای از داده تا تصمیم است. شمارش رخداد به‌تنهایی برای انتخاب محل کلید کافی نیست؛ مدت مواجهه، توپولوژی، محل بار و زمان بازیابی نیز لازم‌اند.

\exercise
اگر دو فیدر هرکدام ۲۰ خاموشی داشته باشند، آیا الزاماً قابلیت اطمینان آن‌ها یکسان است؟
\answer خیر. مدت پوشش داده، زمان هر خاموشی، تعداد مشترکان، مقدار بار و محل رخدادها ممکن است متفاوت باشد؛ بنابراین نرخ خاموشی و \lr{EENS} می‌توانند متفاوت باشند.

\chapter{شناخت و ممیزی داده‌های واقعی کوثر}
\section{فایل‌ها و ساختار داده خاموشی}
پنج فایل اصلی عبارت‌اند از \file{3.xlsx}، \file{4.xlsx}، \file{5.xlsx}، \file{6.xlsx} و \file{7.xlsx}. هرکدام یک شیت با نام دقیق \code{Outages} دارند. یک فایل تک‌ردیفی به نام \file{بدون نام فیدر.xlsx} نیز در آرشیو موجود است. گستره داده‌ها به‌ترتیب در فایل‌های اصلی از \code{A1:O148}، \code{A1:O291}، \code{A1:O343}، \code{A1:O277} و \code{A1:O772} است.

\begin{longtable}{p{0.7cm}p{3.1cm}p{1.6cm}p{2.2cm}p{5.7cm}}
\caption{ستون‌های شیت \code{Outages} با حفظ نام دقیق فایل}\label{tab:columns}\\
\toprule
ردیف & نام ستون & نقش داده‌ای & نوع متداول & تفسیر مهندسی \\
\midrule
\endfirsthead
\toprule ردیف & نام ستون & نقش داده‌ای & نوع متداول & تفسیر مهندسی \\
\midrule
\endhead
۱ & تاخير در ثبت & ثبت اداری & عددی\lr{/}خالی & فاصله ثبت رخداد تا زمان وقوع یا ثبت سامانه‌ای؛ برای نرخ خرابی مستقیماً استفاده نمی‌شود.\\
۲ & تاخير در وصل & بازیابی & عددی\lr{/}خالی & تأخیر ثبت‌شده در فرآیند وصل؛ باید با مدت قطع تطبیق داده شود.\\
۳ & علت قطع & علت خام & متن & مانند «عیب گذرا»، «برخورد پرندگان» یا «بارندگی».\\
۴ & ثبت کننده & ممیزی & متن & نام ثبت‌کننده؛ شاخص فنی خاموشی نیست.\\
۵ & جريان زمان قطع & اندازه الکتریکی & عددی & جریان هنگام قطع؛ مقدار و واحد باید با سامانه اندازه‌گیری کنترل شود.\\
۶ & قطع کننده & تجهیز عمل‌کرده & متن & مانند «بریکر پست فوق توزیع»، «کات اوت ترانس» و «کات اوت سرخط».\\
۷ & توضيحات حادثه & شرح تکمیلی & متن & شرح رویداد، وضعیت فیدر و نحوه وصل مجدد.\\
۸ & انرژي توزيع نشده & شاخص انرژی خام & عددی & واحد در فایل خام صریحاً مشخص نشده است؛ بدون تأیید منبع نباید به‌طور قطعی \lr{kWh} تلقی شود.\\
۹ & نام پست & شناسه پست & متن & انتساب رکورد به پست.\\
۱۰ & نام فيدر & شناسه فیدر & متن\lr{/}عدد & مقدار ثبت‌شده فیدر؛ برای اتصال به فیدر \lr{GIS} باید نرمال‌سازی و ممیزی شود.\\
۱۱ & مدت قطع & زمان خاموشی & زمان اکسل & در فایل‌های نمونه قالب زمان \lr{Excel} دارد؛ تبدیل واحد در بخش ۳.۲ آمده است.\\
۱۲ & تاريخ حادثه & زمان وقوع & تاریخ شمسی & محور اصلی محاسبه مدت پوشش آماری.\\
۱۳ & زمان حادثه & ساعت وقوع & زمان اکسل & برای تحلیل الگوی ساعتی قابل استفاده است.\\
۱۴ & شماره حادثه & شناسه یکتا/عملیاتی & عدد/متن & برای جلوگیری از شمارش تکراری و ردیابی رخداد.\\
۱۵ & انتخاب & پرچم سامانه‌ای & عدد/خالی & ستون عملیاتی؛ به‌تنهایی معنای فنی رخداد را تعیین نمی‌کند.\\
\bottomrule
\end{longtable}

\section{تعداد رکورد و پوشش تاریخ}
\begin{table}[ht]
\centering\small
\caption{مشخصات پایه فایل‌های خاموشی پیوست‌شده}\label{tab:file-meta}
\begin{tabular}{lrrrr}
\toprule
فایل & ردیف داده & نخستین تاریخ & آخرین تاریخ & شیت \\
\midrule
\file{3.xlsx} & 147 & 1400/07/06 & 1405/06/31 & \code{Outages}\\
\file{4.xlsx} & 290 & 1400/07/11 & 1405/06/31 & \code{Outages}\\
\file{5.xlsx} & 342 & 1400/08/15 & 1405/06/31 & \code{Outages}\\
\file{6.xlsx} & 276 & 1400/07/23 & 1405/06/31 & \code{Outages}\\
\file{7.xlsx} & 771 & 1400/07/01 & 1405/06/30 & \code{Outages}\\
\file{بدون نام فیدر.xlsx} & 1 & ناموجود & ناموجود & \code{Outages}\\
\bottomrule
\end{tabular}
\end{table}

وجود یک ردیف در \file{بدون نام فیدر.xlsx} به معنای یک خاموشی قابل انتساب به فیدر مشخص نیست. این رکورد فاقد اطلاعات کلیدی کافی است و باید در فهرست داده‌های نیازمند اصلاح باقی بماند؛ نباید آن را به هیچ‌یک از فیدرهای ۳ تا ۷ نسبت داد.

\section{مطالعه موردی مجزای هر فایل}
\subsection{فایل \file{3.xlsx}؛ فیدر \code{F3}}
این فایل ۱۴۷ رکورد دارد. در طبقه‌بندی خودکار، ۶۲ رخداد در گروه «قابل کاهش با کلیدگذاری»، ۶ رخداد «نیازمند تعمیر یا تعویض تجهیز»، ۴۰ رخداد «ناشی از هماهنگی حفاظتی»، ۵ رخداد «خارج از دامنه» و ۳۴ رخداد «داده ناکافی» قرار گرفته‌اند. بنابراین تعداد رخدادهای تصادفیِ قابل ورود به مدل نرخ خاموشی برابر ۱۰۸ است.

نمونه شماره حادثه \lr{10055783} در تاریخ \lr{1405/04/14} با علت «عیب گذرا» و تجهیز «بریکر پست فوق توزیع» ثبت شده است. شرح «عیب گذرا وصل موفق فیدر بااولین درخواست وصل توسط تکنیسین» و مدت ۵ دقیقه دارد. این نمونه طبق قاعده \Rrule{R-SWITCH-01} در گروه قابل کاهش با کلیدگذاری قرار می‌گیرد. در مقابل، رخداد \lr{3429379} در تاریخ \lr{1400/07/06} علت «برخورد پرندگان» و تجهیز «کات اوت سرخط» دارد؛ چون نقش دقیق تجهیز سرخط می‌تواند روی کل فیدر اثر بگذارد، طبقه‌بند آن را با قاعده \Rrule{R-NAK-01} به بازبینی انسانی می‌فرستد.

\subsection{فایل \file{4.xlsx}؛ فیدر \code{F4}}
این فایل ۲۹۰ رکورد دارد: ۴۴ مورد قابل کاهش با کلیدگذاری، ۶۷ مورد ناشی از هماهنگی حفاظتی، ۱۱ مورد نیازمند تعمیر/تعویض و ۱۶۸ مورد داده ناکافی. تعداد رخدادهای تصادفی واردشده به مدل ۱۲۲ است. نمونه \lr{10047773} در تاریخ \lr{1405/03/24} با علت «عیب گذرا»، تجهیز «بریکر پست فوق توزیع» و مدت ثبت‌شده صفر دقیقه، به گروه قابل کاهش با کلیدگذاری تعلق دارد. صفر دقیقه را نباید خودکار حذف کرد؛ این رخداد می‌تواند عملکرد لحظه‌ای و وصل موفق سریع را نشان دهد.

\subsection{فایل \file{5.xlsx}؛ فیدر \code{F5}}
این فایل ۳۴۲ رکورد دارد: ۱۱۳ مورد قابل کاهش با کلیدگذاری، ۸۶ مورد ناشی از هماهنگی حفاظتی، ۱۶ مورد نیازمند تعمیر/تعویض، ۴ مورد خارج از دامنه و ۱۲۳ مورد داده ناکافی. تعداد رخداد تصادفی ۲۱۵ است.

نمونه \lr{3429550} در تاریخ \lr{1400/09/11} با علت «عیب گذرا» و تجهیز «بریکر پست فوق توزیع» ثبت شده است. شرح حادثه می‌گوید پس از اولین وصل دستی، تمام فیدر به‌صورت متعادل برقدار شد. مدت قطع ۵ دقیقه و مقدار ثبت‌شده در ستون انرژي توزيع نشده برابر ۰٫۰۲۸۶ است. طبقه‌بندی آن \Rrule{R-SWITCH-01} است. در مقابل، رخداد \lr{3429708} با «برخورد پرندگان» و «کات اوت سرخط» با وجود مدت ۲۸ دقیقه، به دلیل ابهام در نقش نقطه‌ای یا سرخطی تجهیز، «داده ناکافی» است.

نمونه \lr{3429483} با علت «عیب گذرا» و شرح «فیدر 5ریکلوز شد» مدت صفر دقیقه دارد. این رخداد برای بحث ریکلوزر و خاموشی‌های گذرا مهم است، ولی صفر بودن مدت به معنای بی‌اهمیت بودن رخداد حفاظتی نیست.

\subsection{فایل \file{6.xlsx}؛ فیدر \code{F6}}
این فایل ۲۷۶ رکورد دارد: ۴۵ مورد قابل کاهش با کلیدگذاری، ۱۱۰ مورد ناشی از هماهنگی حفاظتی، ۱۹ مورد نیازمند تعمیر/تعویض، ۲ مورد خارج از دامنه و ۱۰۰ مورد داده ناکافی. تعداد رخداد تصادفی ۱۷۴ است. نمونه \lr{3429421} در تاریخ \lr{1400/07/23} با علت «عیب گذرا» و شرح «ارت فالت لحظه ای و عملکرد موفق ریکلوزر» دارای مدت صفر دقیقه است و در گروه قابل کاهش با کلیدگذاری قرار دارد. نمونه \lr{3429458} با علت «برخورد پرندگان» و تجهیز «کات اوت ترانس» مدت ۱۱ دقیقه دارد و به علت اثر هماهنگی پایین‌دست، در گروه «ناشی از هماهنگی حفاظتی» قرار می‌گیرد.

\subsection{فایل \file{7.xlsx}؛ فیدر \code{F7}}
این فایل ۷۷۱ رکورد دارد: ۷۰ مورد قابل کاهش با کلیدگذاری، ۱۸۸ مورد ناشی از هماهنگی حفاظتی، ۴۱ مورد نیازمند تعمیر/تعویض، ۷ مورد خارج از دامنه و ۴۶۵ مورد داده ناکافی. تعداد رخداد تصادفی ۲۹۹ است. نمونه \lr{3429389} در تاریخ \lr{1400/07/11}، علت «عیب گذرا»، تجهیز «بریکر پست فوق توزیع» و مدت ۱۸ دقیقه دارد و با قاعده \Rrule{R-SWITCH-01} طبقه‌بندی شده است. رخداد \lr{3429359} با علت «برخورد پرندگان» و «کات اوت ترانس» مدت ۲۱ دقیقه دارد و به گروه «ناشی از هماهنگی حفاظتی» تعلق می‌گیرد.

\section{داده‌های مکانی و فایل \lr{GIS}}
در آرشیو \file{Kosar.zip}، لایه‌های شیپ‌فایل متعددی در مسیر \file{Export\_\allowbreak 1404\_\allowbreak 08\_\allowbreak 17\_\allowbreak 12\_\allowbreak 26\_\allowbreak 05} وجود دارد. برای مطالعه نمونه، دو لایه مهم عبارت‌اند از \code{ELWDTR} برای ترانسفورماتورها و \code{ELWOML} برای خطوط هوایی فشار متوسط. فیلد \code{T\_FeederNa} شناسه فیدر را نگه می‌دارد. در لایه \code{ELWDTR}، فیلد \code{T\_strn} ظرفیت نامی ترانسفورماتور بر حسب \lr{kVA} را ثبت می‌کند؛ در \code{ELWOML}، فیلد \code{Shape\_STLe} طول هندسی ثبت‌شده را نشان می‌دهد.

\begin{table}[ht]
\centering\small\setlength{\tabcolsep}{3pt}
\caption{خلاصه‌ای از ویژگی‌های \lr{GIS} استخراج‌شده برای فیدرهای کوثر}\label{tab:gis}
\begin{tabular}{lrrrr}
\toprule
فیدر & تعداد ترانس در \code{ELWDTR} & جمع \code{T\_strn} (\lr{kVA}) & تعداد قطعات \code{ELWOML} & جمع \code{Shape\_STLe} (\lr{m})\\
\midrule
\code{F3 - KOSAR} & 84 & 7880 & 1332 & 67917.950\\
\code{F5 - KOSAR} & 78 & 6830 & 2043 & 123159.173\\
\code{F6 - KOSAR} & 67 & 8300 & 1189 & 56996.833\\
\code{F7 - KOSAR} & 118 & 9680 & 2519 & 135307.616\\
\bottomrule
\end{tabular}
\end{table}

\warn در داده مکانی، نام فیدر چهارم در بعضی رکوردها به‌صورت \lr{F4- KOSAR} و در نام‌گذاری مورد انتظار به‌صورت \lr{F4 - KOSAR} دیده می‌شود. تفاوت فاصله‌ها می‌تواند باعث عدم اتصال رکوردها در یک پیوند متنی ساده شود. راه درست، نگه‌داشتن مقدار خام و ساخت یک کلید نرمال‌شده جداگانه برای مقایسه است؛ نباید نام اصلی را بی‌ردپا بازنویسی کرد.

\summary{۲}
فایل‌های اصلی هرکدام یک شیت \code{Outages} و ۱۵ ستون مشترک دارند. داده مکانی می‌تواند تعداد و ظرفیت نامی ترانس‌ها و طول هندسی خطوط را بدهد، اما به‌تنهایی بار واقعی لحظه وقوع را تعیین نمی‌کند. فایل بدون نام فیدر قابل تخصیص آماری نیست.

\exercise
در \file{5.xlsx} تعداد رکوردها ۳۴۲ است، اما تعداد رخدادهای تصادفی ۲۱۵ است. کدام عدد باید در صورت کسر نرخ خاموشی تصادفی قرار گیرد؟
\answer ۲۱۵؛ نرخ بر پایه رخدادهای تصادفی معتبر است.

\chapter{پاک‌سازی داده، زمان و پوشش آماری}
\section{چرا پاک‌سازی پیش از مدل‌سازی ضروری است؟}
اگر یک مدت زمان به‌جای دقیقه، به‌صورت کسری از روز \lr{Excel} خوانده شود، خطای مقیاس ۱۴۴۰ برابری رخ می‌دهد. اگر نام یک فیدر به اشتباه به فیدر دیگر متصل شود، نرخ خرابی و بار آن فیدر جابه‌جا خواهد شد. اگر رخداد تکراری دوبار شمرده شود، نرخ خاموشی بیش‌برآورد می‌شود. بنابراین پاک‌سازی صرفاً کار اداری نیست؛ بخشی از مدل ریاضی است.

\section{تبدیل زمان Excel به دقیقه}
در قالب زمانی \lr{Excel}، یک روز برابر عدد ۱ است. بنابراین زمان ذخیره‌شده به‌صورت عدد کسری از روز تبدیل می‌شود. تبدیل عمومی به دقیقه چنین است:
\begin{equation}
 t_{\min}=1440\,d_{\mathrm{Excel}}=60h+m+\frac{s}{60}.
 \label{eq:excel-time}
\end{equation}
در رابطه \eqref{eq:excel-time}، $t_{\min}$ مدت بر حسب دقیقه، $d_{\mathrm{Excel}}$ مقدار کسری از روز در فایل \lr{Excel}، $h$ ساعت، $m$ دقیقه و $s$ ثانیه است. ضریب ۱۴۴۰ برابر تعداد دقیقه‌های شبانه‌روز است.

برای نمونه، در رخداد \lr{3429550} از \file{5.xlsx}، مقدار خام زمانی تقریباً $0.0034722222$ روز است. بنابراین:
\begin{equation}
 t_{\min}=1440\times0.0034722222\simeq5\unit{min}.
 \label{eq:f5-time-example}
\end{equation}
در این مثال، مقدار ورودی سلول، مدت خاموشی محاسبه‌شده و واحد تبدیل‌شده به‌ترتیب $0.0034722222$ روز و ۵ دقیقه‌اند. اگر عدد خام به‌اشتباه ۰٫۰۰۳۴۷ دقیقه خوانده شود، مدت خاموشی تقریباً صفر تلقی خواهد شد.

\section{پاک‌سازی متن و شناسه‌ها}
در متن‌های فارسی ممکن است فاصله معمولی و فاصله نشکن (مانند U+00A0) وجود داشته باشد. بهتر است پیش از مقایسه، فاصله‌های ابتدا و انتها حذف، فاصله‌های چندگانه یکسان و حروف هم‌ارز طبق سیاست مشخص نرمال شوند. با این حال، متن خام باید حفظ شود تا در بازبینی مهندسی بتوان تشخیص داد الگوریتم چه تغییری انجام داده است.

برای شناسه رخداد، اصل کلیدی این است که شماره حادثه \code{شماره حادثه} باید تا حد امکان یکتا باشد. اگر دو ردیف با یک شماره حادثه دیده شوند، پیش از حذف یکی باید معلوم شود که آن‌ها تکرار واقعی‌اند یا دو مرحله از یک رخداد.

\section{پوشش آماری و مواجهه زمانی}
نرخ رخداد باید بر مدت زمانی تقسیم شود که داده در آن معتبر و قابل مشاهده بوده است. فرض کنید بازه بررسی از ماه نخست $M_1$ تا ماه آخر $M_2$ است. هر ماهی که در آن داده معتبر و قابل استفاده وجود دارد، یک ماه پوشش‌دار محسوب می‌شود. ماه‌های بدون مشاهده نباید خودکار به‌عنوان «ماه بدون خاموشی» در نظر گرفته شوند.

\begin{equation}
 T_{\mathrm{exp}}=\sum_{m\in\mathcal{M}_{\mathrm{covered}}} d_m,
 \qquad
 Y_{\mathrm{exp}}=\frac{T_{\mathrm{exp}}}{365}.
 \label{eq:exposure}
\end{equation}
در رابطه \eqref{eq:exposure}، $T_{\mathrm{exp}}$ تعداد روزهای مواجهه معتبر، $\mathcal{M}_{\mathrm{covered}}$ مجموعه ماه‌های دارای پوشش داده، $d_m$ تعداد روزهای ماه $m$ و $Y_{\mathrm{exp}}$ مدت مواجهه بر حسب سال است. در پیاده‌سازی مورد استفاده، روزهای ماه شمسی لحاظ شده و سال مواجهه با تقسیم مجموع روزها بر ۳۶۵ محاسبه می‌شود.

برای \file{5.xlsx}، مدت مواجهه ۱۷۹۶ روز، معادل ۴٫۹۲۰۵۴۸ سال و ۵۹ ماه پوشش‌دار است. در \file{3.xlsx}، ۶۰ ماه تقویمی بین ابتدا و انتهای بازه وجود دارد، اما فقط ۵۱ ماه پوشش‌دار و ۹ ماه فاقد رکوردند؛ در نتیجه مواجهه ۱۵۵۴ روز یا ۴٫۲۵۷۵۳۴ سال است. بنابراین استفاده از فاصله ساده شش‌ساله برای هر پنج فیدر، نرخ‌ها را مخدوش می‌کند.

\section{مدیریت مقدارهای صفر و ناقص}
مدت صفر می‌تواند نشان دهد عملکرد حفاظتی و وصل مجدد در زمان بسیار کوتاه رخ داده است. صفر را نباید بدون بررسی با داده مفقود یکسان گرفت. از طرف دیگر، سلول خالی مدت یا تاریخ یک نقص داده‌ای است، نه یک رخداد صفر دقیقه‌ای. فایل \file{بدون نام فیدر.xlsx} نمونه‌ای است که به دلیل فقدان اطلاعات محوری نباید به یک فیدر مشخص تخصیص داده شود.

\summary{۳}
واحد زمان \lr{Excel} باید پیش از محاسبه به دقیقه تبدیل شود. مخرج نرخ خرابی، مدت مواجهه معتبر است، نه صرفاً فاصله ابتدا تا انتهای فایل. مقدار صفر واقعی و مقدار ناموجود دو وضعیت متفاوت‌اند.

\exercise
اگر مدت \lr{Excel} برابر $0.00625$ روز باشد، مدت آن چند دقیقه است؟
\answer با استفاده از معادله \eqref{eq:excel-time}: $0.00625\times1440=9$ دقیقه.

\chapter{طبقه‌بندی رخداد و مدل قابلیت اطمینان}
\section{شهود فیزیکی طبقه‌بندی}
همه خاموشی‌های ثبت‌شده با یک تجهیز و یک نوع راهکار قابل کاهش نیستند. یک خطای گذرا که بریکر پست عمل کرده و سپس فیدر با موفقیت وصل شده است، با خرابی فیوز یک ترانس پایین‌دست تفاوت دارد. در حالت دوم، ممکن است تعویض فیوز یا رفع خرابی تجهیز لازم باشد و نصب یک کلید در محل دیگر الزاماً مدت خاموشی مصرف‌کننده معیوب را حذف نکند.

در منطق پروژه شش دسته اصلی به کار می‌رود:
\begin{enumerate}
\item \textbf{قابل کاهش با کلیدگذاری:} رخدادهای تصادفی یا گذرا که با جداسازی و مانور مناسب، امکان کاهش خاموشی بار سالم وجود دارد.
\item \textbf{نیازمند تعمیر یا تعویض تجهیز:} رخدادهایی که رفع فیزیکی خرابی تجهیز، شرط اصلی بازیابی است.
\item \textbf{ناشی از هماهنگی حفاظتی:} عملکرد یک تجهیز حفاظتی پایین‌دست که به هماهنگی، تنظیم یا محدوده حفاظتی وابسته است.
\item \textbf{برنامه‌ریزی‌شده:} خاموشی‌های ناشی از عملیات برنامه‌ریزی‌شده، در صورتی که از متن قابل تشخیص باشند.
\item \textbf{خارج از دامنه:} رویدادهایی که طبق محدوده مطالعه نباید وارد مدل شوند.
\item \textbf{داده ناکافی:} رخداد فاقد اطلاعات کافی یا دارای ابهام مهم در محل و نقش تجهیز.
\end{enumerate}

\section{قواعد نمونه و معنای کدها}
در نسخه بررسی‌شده، کدهای کلیدی شامل موارد زیرند: \Rrule{R-NAK-01} برای داده ناکافی، \Rrule{R-PLAN-01} برای برنامه‌ریزی‌شده، \Rrule{R-UNK-99} برای علت نگاشت‌نشده و \Rrule{R-DOM-01} برای خارج از دامنه. همچنین \Rrule{R-REP-01} برای تعمیر/تعویض تجهیز، \Rrule{R-COORD-01} برای هماهنگی حفاظتی و \Rrule{R-SWITCH-01} برای رخداد قابل کاهش با کلیدگذاری به کار می‌روند.

قاعده مهم این است که «برخورد پرندگان» به‌تنهایی برای قرار دادن رخداد در گروه قابل کاهش کافی نیست. در \file{3.xlsx}، رخداد \lr{3429379} با همین علت اما تجهیز «کات اوت سرخط» در گروه داده ناکافی قرار گرفته است؛ در مقابل، رخداد \lr{10055783} با «عیب گذرا» و «بریکر پست فوق توزیع» قابل کاهش با کلیدگذاری است. محل و نقش تجهیز، در کنار علت خام، تعیین‌کننده‌اند.

\section{نرخ رخداد تصادفی و نرخ قابل کاهش}
برای محاسبه نرخ، ابتدا رخدادهای تصادفی را از رخدادهای برنامه‌ریزی‌شده، خارج از دامنه و داده ناکافی جدا می‌کنیم. در منطق فعلی، دسته‌های ۱، ۲ و ۳ در شمارش رخداد تصادفی قرار می‌گیرند؛ اما فقط دسته ۱ «قابل کاهش با کلیدگذاری» است. این دو شمارش را نباید با یکدیگر اشتباه گرفت.

\begin{equation}
 \lambda=\frac{N_{\mathrm{random}}}{Y_{\mathrm{exp}}},
 \qquad
 \lambda_{\mathrm{red}}=\frac{N_{\mathrm{reducible}}}{Y_{\mathrm{exp}}}.
 \label{eq:lambda}
\end{equation}
در معادله \eqref{eq:lambda}، $\lambda$ نرخ سالانه رخدادهای تصادفی، $\lambda_{\mathrm{red}}$ نرخ سالانه رخدادهای قابل کاهش با کلیدگذاری، $N_{\mathrm{random}}$ تعداد رخدادهای تصادفی، $N_{\mathrm{reducible}}$ تعداد رخدادهای قابل کاهش و $Y_{\mathrm{exp}}$ مدت مواجهه بر حسب سال است. واحد هر نرخ، رخداد بر سال است.

برای فیدر \code{F5} از \file{5.xlsx}، تعداد رخداد تصادفی ۲۱۵، رخداد قابل کاهش ۱۱۳ و مواجهه ۴٫۹۲۰۵۴۸ سال است:
\begin{align}
 \lambda_{F5}&=\frac{215}{4.920548}=43.6943\;\mathrm{event/year},\\
 \lambda_{\mathrm{red},F5}&=\frac{113}{4.920548}=22.9649\;\mathrm{event/year}.
 \label{eq:f5-lambda}
\end{align}
در این محاسبه، صورت کسر اول تمام رخدادهای تصادفی سه دسته اول و صورت کسر دوم فقط دسته «قابل کاهش با کلیدگذاری» است.

\section{مدت متوسط خاموشی و دسترس‌پذیری}
مدت متوسط خاموشی رخدادهای تصادفی از تقسیم مجموع مدت رخدادها بر تعداد همان رخدادها به دست می‌آید:
\begin{equation}
 r=\frac{\sum_{j=1}^{N_{\mathrm{random}}}t_j}{N_{\mathrm{random}}}.
 \label{eq:r}
\end{equation}
در رابطه \eqref{eq:r}، $r$ مدت متوسط خاموشی بر حسب دقیقه، $t_j$ مدت رخداد $j$ و $N_{\mathrm{random}}$ تعداد رخدادهای تصادفی است. اگر رخدادهای صفر دقیقه‌ای معتبر باشند، در مخرج باقی می‌مانند.

شاخص عدم‌دسترس‌پذیری مورد انتظار بر حسب ساعت در سال، با فرض همگن بودن رخدادها، به‌صورت زیر محاسبه می‌شود:
\begin{equation}
 U=\frac{\lambda r}{60}.
 \label{eq:unavailability}
\end{equation}
در رابطه \eqref{eq:unavailability}، $U$ ساعات خاموشی مورد انتظار در سال، $\lambda$ نرخ رخداد بر سال، $r$ مدت متوسط هر رخداد بر دقیقه و عدد ۶۰ ضریب تبدیل دقیقه به ساعت است.

برای \code{F5}، مقدار $r=19.051$ دقیقه و $\lambda=43.6943$ رخداد بر سال است؛ بنابراین $U\simeq13.8737$ ساعت در سال. این مقدار یک شاخص برآوردی مبتنی بر داده و فرض نرخ پایا است، نه برنامه زمانی خاموشی آینده.

\section{شاخص‌های مبتنی بر مشترکین: \lr{SAIFI}، \lr{SAIDI}، \lr{CAIDI} و \lr{ASAI}}
در ارزیابی شرکت توزیع، نرخ خاموشی فیدر به‌تنهایی کافی نیست؛ تعداد مشترکینی که در هر رخداد تحت تأثیر قرار می‌گیرند نیز اهمیت دارد. شاخص‌های متداول به‌صورت زیر تعریف می‌شوند:
\begin{align}
 SAIFI&=\frac{\sum_i\lambda_i N_i}{N_T},\label{eq:saifi}\\
 SAIDI&=\frac{\sum_i U_i N_i}{N_T},\label{eq:saidi}\\
 CAIDI&=\frac{SAIDI}{SAIFI},\label{eq:caidi}\\
 ASAI&=1-\frac{SAIDI}{8760}.\label{eq:asai}
\end{align}
در معادلات \eqref{eq:saifi} تا \eqref{eq:asai}، $SAIFI$ متوسط تعداد وقفه به‌ازای هر مشترک در سال، $SAIDI$ متوسط مدت قطع به‌ازای هر مشترک بر حسب ساعت در سال، $CAIDI$ متوسط مدت هر وقفه به ازای هر مشترک بر حسب ساعت، $ASAI$ شاخص متوسط دسترس‌پذیری، $N_i$ تعداد مشترکین متأثر از رخداد/بخش $i$، $N_T$ کل مشترکین تحت بررسی، $\lambda_i$ نرخ رخداد بخش $i$ و $U_i$ مدت قطع سالانه بخش $i$ بر حسب ساعت است. عدد ۸۷۶۰ تعداد ساعت‌های یک سال غیرکبیسه است.

\warn فایل‌های خاموشی کوثر ستون صریح تعداد مشترکین قطع‌شده و کل مشترکین فیدر را ندارند. از این رو، $SAIFI$، $SAIDI$، $CAIDI$ و $ASAI$ را نمی‌توان فقط با داده‌های پیوست‌شده به‌صورت معتبر عددگذاری کرد. برای نمونه، رخداد ۵ دقیقه‌ای \lr{3429550} در \file{5.xlsx} مدت را مشخص می‌کند، اما تعداد مشترکین متأثر را تعیین نمی‌کند.

\section{نتایج مقایسه‌ای پنج فیدر}
\begin{longtable}{lrrrrrrrr}
\caption{نتایج تحلیل رخدادها در فایل‌های خاموشی پیوست‌شده}\label{tab:reliability}\\
\toprule
فایل & کل & تصادفی & قابل کاهش & داده ناکافی & هماهنگی & تعمیر & $\lambda$ & $r$ (دقیقه)\\
\midrule
\endfirsthead
\toprule فایل & کل & تصادفی & قابل کاهش & داده ناکافی & هماهنگی & تعمیر & $\lambda$ & $r$ (دقیقه)\\
\midrule
\endhead
\file{3.xlsx} &147&108&62&34&40&6&25.367&24.315\\
\file{4.xlsx} &290&122&44&168&67&11&26.102&32.148\\
\file{5.xlsx} &342&215&113&123&86&16&43.694&19.051\\
\file{6.xlsx} &276&174&45&100&110&19&37.249&20.184\\
\file{7.xlsx} &771&299&70&465&188&41&59.767&36.214\\
\bottomrule
\end{longtable}
\begin{table}[ht]
\centering\small
\caption{نرخ قابل کاهش، مدت متوسط رخدادهای قابل کاهش و عدم‌دسترس‌پذیری سالانه}\label{tab:lambda-u}
\begin{tabular}{lrrrrr}
\toprule فایل & مواجهه (سال) & $\lambda_{\mathrm{red}}$ & $r_{\mathrm{red}}$ (دقیقه) & $U$ (ساعت/سال) & پوشش/مفقود\\
\midrule
\file{3.xlsx}&4.257534&14.562&21.484&10.2799&51/9\\
\file{4.xlsx}&4.673973&9.414&10.182&13.9854&56/4\\
\file{5.xlsx}&4.920548&22.965&2.841&13.8737&59/0\\
\file{6.xlsx}&4.671233&9.633&9.067&12.5307&56/5\\
\file{7.xlsx}&5.002740&13.992&13.814&36.0735&60/0\\
\bottomrule
\end{tabular}
\end{table}

\section{تفسیر مهندسی نتایج}
فیدر \code{F7} بیشترین نرخ رخداد تصادفی را در این مجموعه دارد، اما ۴۶۵ رکورد آن داده ناکافی‌اند؛ بنابراین نتیجه‌گیری درباره عملکرد فیدر باید همراه با هشدار کیفیت داده باشد. فیدر \code{F5} نرخ قابل کاهش بالاتری نسبت به دیگر فیدرها دارد (۲۲٫۹۶۵ رخداد در سال) و از این جهت مطالعه جالبی برای سنجش فایده کلیدگذاری است. با این حال، نرخ بالا به‌تنهایی محل نصب کلید را تعیین نمی‌کند؛ بار پایین‌دست و ساختار شبکه نیز لازم‌اند.

\warn مقدار \code{energy\_kwh} در خروجی محاسباتی پروژه از ستون خام انرژی نگاشت شده، اما مستندات بررسی‌شده تصریح می‌کنند که واحد منبع صریحاً تعریف نشده است. بنابراین جمع‌های این ستون در این آموزش به‌عنوان انرژی قطعی \lr{kWh} برای محاسبه \lr{ENS} استفاده نمی‌شوند.

\summary{۴}
نرخ خاموشی باید با مدت مواجهه محاسبه شود. در مدل داده حاضر، نرخ تصادفی و نرخ قابل کاهش دو مفهوم جدا هستند. برای \code{F5}، به‌ترتیب ۴۳٫۶۹۴ و ۲۲٫۹۶۵ رخداد بر سال به دست آمد.

\exercise
برای فیدر \code{F6}، ۱۷۴ رخداد تصادفی و ۴٫۶۷۱۲۳۳ سال مواجهه داریم. نرخ رخداد تقریباً چقدر است؟
\answer $\lambda=174/4.671233\simeq37.249$ رخداد بر سال.

\chapter{از خاموشی تا انرژی تأمین‌نشده: \lr{ENS} و \lr{EENS}}
\section{شهود \lr{ENS} و \lr{EENS}}
اگر بار قطع‌شده $P$ کیلووات و مدت قطع $t$ ساعت باشد، انرژی تأمین‌نشده تقریباً برابر حاصل‌ضرب این دو است. \lr{ENS} برای یک رخداد یا یک بازه مشاهده‌شده به کار می‌رود. \lr{EENS} مقدار مورد انتظار انرژی تأمین‌نشده در یک بازه مرجع، معمولاً یک سال، است. \lr{EENS} علاوه بر توان بار، به نرخ رخداد و زمان بازیابی وابسته است.

\section{فرمول انرژی یک رخداد}
\begin{equation}
 ENS_j=\int_{t_{0,j}}^{t_{1,j}}P_{\mathrm{lost},j}(t)\,dt.
 \label{eq:ens-integral}
\end{equation}
در رابطه \eqref{eq:ens-integral}، $ENS_j$ انرژی تأمین‌نشده رخداد $j$ بر حسب \lr{kWh}، $P_{\mathrm{lost},j}(t)$ توان قطع‌شده در زمان $t$ بر حسب \lr{kW} و $t_{0,j}$ و $t_{1,j}$ زمان شروع و پایان خاموشی بر حسب ساعت‌اند. اگر توان قطع‌شده در مدت خاموشی ثابت باشد:
\begin{equation}
 ENS_j=P_j\,t_j.
 \label{eq:ens-constant}
\end{equation}
در رابطه \eqref{eq:ens-constant}، $P_j$ توان قطع‌شده ثابت بر حسب \lr{kW} و $t_j$ مدت قطع بر حسب ساعت است.

\subsection{مثال واقعی رخداد با بار فرضی شفاف}
برای رخداد \lr{3429550} در \file{5.xlsx}، مدت خاموشی ۵ دقیقه است. در داده‌های \lr{GIS}، جمع ظرفیت نامی ترانس‌های فیدر \code{F5} در لایه \code{ELWDTR} برابر ۶۸۳۰ \lr{kVA} است. ظرفیت نامی ترانس، بار واقعی لحظه‌ای نیست. برای مثال آموزشی فرض می‌کنیم ضریب بار $LF=0.50$ و ضریب توان $PF=0.90$ باشد. در این صورت:
\begin{equation}
 P_{\mathrm{est}}=S_{\mathrm{rated}}\times LF\times PF
 =6830\times0.50\times0.90=3073.5\unit{kW}.
 \label{eq:f5-load-est}
\end{equation}
در معادله \eqref{eq:f5-load-est}، $S_{\mathrm{rated}}$ مجموع ظرفیت نامی ترانس‌ها بر حسب \lr{kVA}، $LF$ ضریب بار، $PF$ ضریب توان و $P_{\mathrm{est}}$ برآورد آموزشی بار اکتیو بر حسب \lr{kW} است. مقدار ۳۰۷۳٫۵ \lr{kW} نتیجه فرض‌های $LF$ و $PF$ است و اندازه‌گیری بار فیدر نیست.

اگر برای یک برآورد ساده فرض کنیم کل بار برآوردی در آن پنج دقیقه قطع بوده است:
\begin{equation}
 ENS_{3429550}=3073.5\times\frac{5}{60}=256.125\unit{kWh}.
 \label{eq:f5-ens}
\end{equation}
در رابطه \eqref{eq:f5-ens}، توان برآوردی ۳۰۷۳٫۵ \lr{kW} و مدت قطع ۵ دقیقه است که به ساعت تبدیل شده است. مقدار ۲۵۶٫۱۲۵ \lr{kWh} صرفاً خروجی سناریوی آموزشی با فرض قطع کل بار است؛ \lr{ENS} واقعی نیازمند بار واقعی لحظه رخداد و مقدار بار واقعاً قطع‌شده است.

\section{\lr{EENS} برای بخش‌های فیدر}
برای هر بخش $i$، فرض می‌کنیم رخداد با نرخ $\lambda_i$ در سال رخ می‌دهد. در زمان تعمیر، بار بخش معیوب $P_i$ قطع است. بخشی از بار سالم ممکن است پس از مانور در زمان $T_{\mathrm{restore}}$ برقدار شود؛ بار سالمی که تا تعمیر نهایی بدون برق می‌ماند نیز در محاسبه می‌آید.

\begin{equation}
 EENS_i=\frac{\lambda_i}{60}\left[T_{\mathrm{rep}}P_i+T_{\mathrm{rest}}P_{\mathrm{restored}}+T_{\mathrm{rep}}P_{\mathrm{remaining}}\right].
 \label{eq:eens-section}
\end{equation}
در رابطه \eqref{eq:eens-section}، $EENS_i$ انرژی مورد انتظار سالانه ناشی از خطای بخش $i$ بر حسب \lr{kWh/year}، $\lambda_i$ نرخ خطای بخش بر سال، $T_{\mathrm{rep}}$ زمان تعمیر بر حسب دقیقه، $P_i$ بار بخش معیوب بر حسب \lr{kW}، $T_{\mathrm{rest}}$ زمان بازیابی بار سالم بر حسب دقیقه، $P_{\mathrm{restored}}$ بار سالمی که با مانور بازیابی می‌شود و $P_{\mathrm{remaining}}$ بار سالمی است که تا پایان تعمیر بی‌برق می‌ماند. تقسیم بر ۶۰ تبدیل \lr{kW-minute} به \lr{kWh} است.

\begin{equation}
 P_{\mathrm{restored}}=\min(P_{\mathrm{healthy}},C_{\mathrm{backup}}),
 \qquad
 P_{\mathrm{remaining}}=P_{\mathrm{healthy}}-P_{\mathrm{restored}}.
 \label{eq:backup-cap}
\end{equation}
در رابطه \eqref{eq:backup-cap}، $P_{\mathrm{healthy}}$ مجموع بار سالم قابل انتقال، $C_{\mathrm{backup}}$ ظرفیت انتقال یا پشتیبانی مسیر جایگزین بر حسب \lr{kW} و تابع $\min$ کوچک‌ترین دو مقدار را انتخاب می‌کند. اگر ظرفیت پشتیبان در مدل نامحدود فرض شود، $P_{\mathrm{restored}}=P_{\mathrm{healthy}}$؛ در غیر این صورت ظرفیت محدود، مقدار بازیابی‌شده را محدود می‌کند.

\section{مدل یکنواخت و مدل جهت‌دار}
در مدل \textbf{\lr{uniform}}، فرض می‌شود تمام بار سالمی که از نظر ظرفیت پشتیبان قابل انتقال است، امکان بازیابی دارد؛ جهت فیزیکی مسیر و محل کلید در این ساده‌سازی وارد جزئیات نمی‌شود. در مدل \textbf{\lr{directional}}، فقط بار پایین‌دستِ قابل جداسازی نسبت به تجهیز یا محل خطا واجد بازیابی در نظر گرفته می‌شود. این مدل به ساختار شعاعی و جهت واقعی جریان و مانور نزدیک‌تر است.

در مدل جهت‌دار، بار سالم قابل بازیابی به بار پایین‌دستِ همان بخش محدود می‌شود:
\begin{equation}
 EENS_i^{\mathrm{dir}}=\frac{\lambda_i}{60}\left[T_{\mathrm{rep}}P_i+T_{\mathrm{rest}}P_{\mathrm{down,rest},i}+T_{\mathrm{rep}}(P_{\mathrm{down},i}-P_{\mathrm{down,rest},i})\right].
 \label{eq:eens-directional}
\end{equation}
در رابطه \eqref{eq:eens-directional}، $P_{\mathrm{down},i}$ مجموع بار زون‌های پایین‌دست بخش $i$ و $P_{\mathrm{down,rest},i}$ بار پایین‌دستی‌ای است که با مانور و در محدودیت ظرفیت پشتیبان بازیابی می‌شود. سایر نمادها مطابق معادله \eqref{eq:eens-section} تعریف شده‌اند. این رابطه روشن می‌کند که مدل جهت‌دار از اطلاعات ساختار زنجیره‌ای زون‌ها استفاده می‌کند.

زمان بازیابی قابل استفاده در مدل، از زمان تعمیر بیشتر نمی‌شود:
\begin{equation}
 T_{\mathrm{rest}}=\min(T_{\mathrm{restore}},T_{\mathrm{rep}}).
 \label{eq:trest}
\end{equation}
در رابطه \eqref{eq:trest}، $T_{\mathrm{restore}}$ زمان لازم برای مانور و برقدار کردن بار سالم و $T_{\mathrm{rep}}$ زمان تعمیر تجهیز معیوب است؛ هر دو بر حسب دقیقه‌اند.

\section{جمع \lr{EENS} و درصد بهبود}
\begin{equation}
 EENS_{\mathrm{total}}=\sum_{i=1}^{n}EENS_i,
 \qquad
 \Delta EENS(\%)=100\frac{EENS_{\mathrm{base}}-EENS_{\mathrm{new}}}{EENS_{\mathrm{base}}}.
 \label{eq:eens-total}
\end{equation}
در رابطه \eqref{eq:eens-total}، $n$ تعداد بخش‌های مدل، $EENS_{\mathrm{base}}$ مقدار پیش از نصب یا تغییر تجهیزات و $EENS_{\mathrm{new}}$ مقدار پس از طرح پیشنهادی است. $\Delta EENS$ درصد کاهش انرژی تأمین‌نشده است. اگر مقدار مبنا صفر باشد، این نسبت تعریف‌نشده است.

\section{مثال دو ناحیه‌ای آموزشی بر مبنای فیدر F5}
\assumption برای نشان دادن سازوکار \lr{EENS}، فیدر \code{F5} را به دو ناحیه مساوی از نظر طول و بار فرض می‌کنیم. طول کل ثبت‌شده در \code{ELWOML} برابر ۱۲۳٫۱۵۹ \lr{km} است؛ بنابراین هر نیمه فرضاً ۶۱٫۵۸۰ \lr{km} خواهد بود. بار برآوردی ۳۰۷۳٫۵ \lr{kW} از معادله \eqref{eq:f5-load-est} به دو بار ۱۵۳۶٫۷۵ \lr{kW} تقسیم می‌شود. فرض‌های زمان طراحی $T_{\mathrm{rep}}=50$ دقیقه، $T_{\mathrm{restore}}=20$ دقیقه و ظرفیت پشتیبان نامحدودند. این تقسیم‌بندی یک سناریوی محاسباتی آموزشی است، نه خروجی جایابی واقعی شبکه کوثر.

در حالت بدون جداسازی مؤثر، تمام بار برای ۵۰ دقیقه در معرض قطع است:
\begin{equation}
 EENS_{\mathrm{base}}=\frac{43.6943}{60}\times50\times3073.5
 \simeq111912.1\unit{kWh/year}.
 \label{eq:f5-base-eens}
\end{equation}
در رابطه \eqref{eq:f5-base-eens}، نرخ رخداد تصادفی فیدر \code{F5} و بار برآوردی آن از داده و فرض‌های قبلی گرفته شده‌اند؛ ۵۰ دقیقه زمان تعمیر فرض‌شده است.

پس از جداسازی فرضی دو ناحیه، نرخ هر ناحیه به دلیل فرض طول برابر نصف نرخ کل می‌شود. در هر رخداد، بار بخش معیوب ۵۰ دقیقه و بار سالم ۲۰ دقیقه قطع می‌ماند:
\begin{equation}
 EENS_{\mathrm{new}}=2\times\frac{43.6943/2}{60}\left[50(1536.75)+20(1536.75)\right]
 \simeq78338.5\unit{kWh/year}.
 \label{eq:f5-new-eens}
\end{equation}
در رابطه \eqref{eq:f5-new-eens}، ضریب ۲ دو ناحیه متقارن، $43.6943/2$ نرخ تخصیص‌یافته به هر ناحیه و ۱۵۳۶٫۷۵ \lr{kW} بار هر ناحیه است. کاهش نسبی در این سناریو ۳۰ درصد است. این نتیجه صرفاً برای نمایش اثر بازیابی بار سالم است؛ مقدار واقعی به توپولوژی، مکان کلید، بار واقعی هر زون و محدودیت ظرفیت مسیر جایگزین وابسته است.

\summary{۵}
\lr{ENS} انرژی قطع‌شده در یک رخداد و \lr{EENS} مقدار مورد انتظار سالانه آن است. \lr{EENS} باید بار، نرخ خطا و مدت‌های بازیابی را هم‌زمان در نظر بگیرد. ظرفیت نامی ترانس به‌تنهایی بار واقعی نیست؛ هر تبدیل آن به \lr{kW} نیازمند فرض یا اندازه‌گیری است.

\exercise
اگر بار ثابت ۱۲۰۰ \lr{kW} به مدت ۱۵ دقیقه قطع شود، \lr{ENS} چقدر است؟
\answer $ENS=1200\times15/60=300$ \lr{kWh}.

\chapter{توپولوژی شبکه، تخصیص نرخ و بار زون‌ها}
\section{نمایش شبکه به‌صورت گراف}
برای تحلیل جایابی، شبکه باید به شکل یک گراف مدل شود. گراف از مجموعه رأس‌ها و یال‌ها تشکیل می‌شود. رأس می‌تواند نقطه انشعاب، محل ترانس، کلید یا انتهای خط باشد و یال نماینده یک قطعه خط بین دو رأس است.
\begin{equation}
 G=(V,E),
 \qquad
 e_{uv}=(u,v)\in E.
 \label{eq:graph}
\end{equation}
در معادله \eqref{eq:graph}، $G$ گراف شبکه، $V$ مجموعه رأس‌ها، $E$ مجموعه یال‌ها و $e_{uv}$ یال بین رأس‌های $u$ و $v$ است. در فیدر شعاعی، مسیر یکتایی از منبع تا هر بار وجود دارد.

\section{تقسیم فیدر به زون}
زون یا بخش، قسمتی از فیدر است که با تجهیزات کلیدزنی و نقاط جداسازی تعریف می‌شود. پس از وقوع خطا در یک زون، بار داخل زون معیوب تا تعمیر قطع می‌ماند؛ اما بار سالم ممکن است با باز کردن و بستن کلیدها از سمت دیگر تغذیه شود. تعیین زون‌ها باید با کلیدهای موجود، نقاط انشعاب و امکان مانور سازگار باشد.

\section{تخصیص نرخ خطا بر اساس طول}
اگر اطلاعات مکانیِ نرخ خطای مجزا برای هر قطعه وجود نداشته باشد، یک تقریب متداول این است که نرخ کل فیدر متناسب با طول خطوط بین بخش‌ها توزیع شود:
\begin{equation}
 \lambda_i=\lambda_{\mathrm{feeder}}\frac{L_i}{\sum_{k=1}^{n}L_k},
 \qquad
 \sum_{i=1}^{n}\lambda_i=\lambda_{\mathrm{feeder}}.
 \label{eq:length-lambda}
\end{equation}
در معادله \eqref{eq:length-lambda}، $\lambda_i$ نرخ خطای تخصیص‌یافته به بخش $i$، $\lambda_{\mathrm{feeder}}$ نرخ کل فیدر، $L_i$ طول بخش $i$، $n$ تعداد بخش‌ها و $k$ اندیس جمع است. این روش فرض می‌کند تراکم خرابی به ازای واحد طول در همه قطعات یکسان است.

برای \code{F5}، مجموع طول فیلد \code{Shape\_STLe} در لایه \code{ELWOML} برابر ۱۲۳۱۵۹٫۱۷۳ متر است. در سناریوی دو نیمه مساوی، هر نیمه ۶۱۵۷۹٫۵۸۷ متر و ۵۰ درصد نرخ فیدر را می‌گیرد. با نرخ ۴۳٫۶۹۴۳ رخداد در سال، نرخ هر نیمه ۲۱٫۸۴۷۲ رخداد در سال خواهد بود. این تخصیص یک تقریب طول‌محور است و در نقاطی که نوع هادی، سن خط یا محیط متفاوت است، می‌تواند دقت کافی نداشته باشد.

\section{توزیع بار}
بار هر زون بهتر است از داده اندازه‌گیری، بار ترانس‌ها، منحنی بار یا مدل مشترکین به دست آید. جمع ظرفیت نامی ترانس فقط یک شاخص ظرفیت نصب‌شده است. اگر بار واقعی در دسترس نباشد، می‌توان سناریو تعریف کرد، اما باید ضریب بار و ضریب توان را صریح نوشت و تحلیل حساسیت انجام داد.

\begin{equation}
 P_i=\sum_{j\in\mathcal{T}_i}S_j\,LF_j\,PF_j.
 \label{eq:zone-load}
\end{equation}
در رابطه \eqref{eq:zone-load}، $P_i$ بار اکتیو برآوردی زون $i$، $\mathcal{T}_i$ مجموعه ترانس‌های واقع در آن زون، $S_j$ ظرفیت نامی ترانس $j$ بر حسب \lr{kVA}، $LF_j$ ضریب بار آن ترانس و $PF_j$ ضریب توان آن است. جمع حاصل بر حسب \lr{kW} خواهد بود.

\section{داده \lr{GIS} و محدودیت اتصال}
لایه \code{ELWDTR} برای نسبت‌دادن ترانس به فیدر و استخراج ظرفیت نامی مفید است. لایه \code{ELWOML} هندسه خطوط هوایی را در اختیار می‌گذارد. در یک فرآیند کامل، باید با استفاده از مختصات، اتصال هندسی و توپولوژیک خطوط نیز بازسازی شود؛ نزدیک بودن دو قطعه روی نقشه به‌تنهایی به معنی اتصال الکتریکی نیست. انشعاب، تقاطع بدون اتصال، جامپر و وضعیت باز/بسته کلیدها باید لحاظ شود.

\summary{۶}
مدل گرافی رابطه بین توپولوژی و بار را ساختارمند می‌کند. تخصیص نرخ بر اساس طول یک فرض یکنواختی دارد. \lr{GIS} می‌تواند طول و ظرفیت ترانس را بدهد، اما برای مدل قابل اتکا باید اتصال الکتریکی و بار واقعی نیز اعتبارسنجی شوند.

\exercise
اگر نرخ کل فیدر ۳۰ رخداد بر سال و طول یک زون ۲۰ درصد طول کل باشد، نرخ تخصیص‌یافته به آن زون طبق مدل طول‌محور چقدر است؟
\answer $\lambda_i=30\times0.20=6$ رخداد بر سال.

\chapter{تجهیزات حفاظتی و نقش عملکردی آن‌ها}
\section{فیوز و کات اوت}
فیوز یا کات اوت معمولاً برای حفاظت انشعاب یا ترانس به کار می‌رود. اگر خطا در پایین‌دست فیوز رخ دهد، عملکرد صحیح فیوز می‌تواند محدوده خاموشی را به همان انشعاب محدود کند. اگر هماهنگی حفاظتی مناسب نباشد، ممکن است تجهیز بالادست نیز عمل کند و بخش بزرگ‌تری قطع شود. بنابراین در داده‌های خاموشی، ثبت نام تجهیز عمل‌کرده باید همراه با محل فیزیکی آن بررسی شود.

\section{سکسیونر و سکسیونر قابل قطع زیر بار}
سکسیونر برای ایجاد جدایی قابل رؤیت و ایمن‌سازی بخش شبکه استفاده می‌شود و به‌طور معمول برای قطع جریان خطا طراحی نشده است. سکسیونر قابل قطع زیر بار می‌تواند جریان بار را در محدوده نامی خود قطع کند، اما جایگزین بریکر یا ریکلوزر برای قطع جریان اتصال کوتاه نیست. هنگام طراحی، ظرفیت نامی، قدرت قطع، روش بهره‌برداری و الزامات ایمنی باید از مشخصات سازنده و استاندارد معتبر اخذ شوند.

\section{ریکلوزر}
ریکلوزر یک کلید حفاظتی خودکار است که پس از قطع جریان خطا می‌تواند طبق برنامه، مجدداً وصل کند. کاربرد شاخص آن در خطاهای گذراست. اگر خطا با وصل مجدد رفع شده باشد، بخش بزرگی از بار ممکن است تنها یک وقفه کوتاه را تجربه کند. اگر خطا باقی بماند، ریکلوزر پس از تعداد عملیات تنظیم‌شده به قفل نهایی می‌رود.

در فایل‌های کوثر، شرح‌هایی مانند «فیدر 5ریکلوز شد» در رخداد \lr{3429483} از \file{5.xlsx} و «ارت فالت لحظه ای و عملکرد موفق ریکلوزر» در رخداد \lr{3429421} از \file{6.xlsx} وجود دارد. این شواهد نشان می‌دهد که رخدادهای گذرا باید جدا از خاموشی‌های دائمی تحلیل شوند.

\section{سکشنالایزر}
سکشنالایزر معمولاً با شمارش جریان خطا یا هماهنگی با تجهیز بالادست، پس از بی‌برق شدن مدار باز می‌شود تا بخش معیوب جدا گردد. سکشنالایزر به‌تنهایی جریان خطا را قطع نمی‌کند؛ باید با ریکلوزر یا بریکر بالادست هماهنگ شود. اگر تعداد شمارش، زمان بازشدن یا محل آن با چرخه ریکلوزر سازگار نباشد، ممکن است جداسازی درست انجام نشود.

\section{انتخاب تجهیز بر مبنای نوع رخداد}
\begin{table}[ht]
\centering\small
\caption{ارتباط مفهومی نوع رخداد با ابزار اصلاحی}\label{tab:device-choice}
\begin{tabularx}{\textwidth}{p{3.1cm}p{3.4cm}X}
\toprule
نوع رخداد یا مسئله & راهکار محتمل & شرط بررسی \\
\midrule
خطای گذرا با وصل موفق & ریکلوزر و تنظیم مناسب & نرخ خطاهای گذرا، تعداد شات، زمان مرده و رفتار بار حساس.\\
خطای دائمی روی شاخه & سکشنالایزر یا کلید جداساز همراه با حفاظت بالادست & امکان قطع جریان خطا توسط تجهیز بالادست و امکان ایزوله‌کردن شاخه.\\
خاموشی گسترده ناشی از خطای موضعی & کلیدهای مرزی/سکسیونر و مانور & وجود مسیر جایگزین، ظرفیت پشتیبان و زمان بازیابی.\\
عملکرد نامناسب فیوز پایین‌دست & بازبینی هماهنگی حفاظتی و انتخاب فیوز & منحنی جریان-زمان، جریان اتصال کوتاه و جریان هجومی.\\
خرابی خود تجهیز & تعمیر یا تعویض تجهیز & کلید جدید لزوماً علت اصلی خرابی را برطرف نمی‌کند.\\
\bottomrule
\end{tabularx}
\end{table}

\summary{۷}
هر تجهیز وظیفه مشخصی دارد. ریکلوزر برای قطع و وصل خودکار و خطاهای گذرا مهم است؛ سکشنالایزر برای جداسازی بخش معیوب با هماهنگی تجهیز بالادست است؛ سکسیونر وسیله جداسازی و ایمنی است و نباید به‌عنوان بریکر اتصال کوتاه در نظر گرفته شود.

\exercise
آیا سکشنالایزر می‌تواند به‌تنهایی جریان اتصال کوتاه را قطع کند؟
\answer خیر. سکشنالایزر معمولاً پس از قطع جریان توسط تجهیز بالادست و در هماهنگی با آن باز می‌شود.

\chapter{مدل خطای گذرا و عملکرد ریکلوزر}
\section{چرا مدل گذرا از مدل خطای دائمی جداست؟}
در خطای گذرا، ممکن است علت خطا با قطع موقت جریان برطرف شود؛ در این صورت، وصل مجدد موفق ریکلوزر می‌تواند از باقی‌ماندن خاموشی طولانی جلوگیری کند. اما در خطای دائمی، وصل مجدد ممکن است به قطع مجدد و نهایتاً قفل ریکلوزر بینجامد. بنابراین نمی‌توان همه رخدادها را با یک زمان تعمیر ثابت مدل کرد.

\section{میانگین زمان وقفه بار محافظت‌شده}
برای یک نمایش ساده، فرض کنیم ریکلوزر با احتمال $p$ در وصل مجدد موفق باشد. در مدل نمونه پروژه، مدت وقفه گذرا $t_{\mathrm{tr}}$، زمان اتصال مجدد $t_{\mathrm{reclose}}$ و احتمال موفقیت $p$ هستند. زمان مورد انتظار وقفه بار تحت حفاظت چنین نوشته می‌شود:
\begin{equation}
 t_{\mathrm{prot,eff}}=t_{\mathrm{reclose}}+(1-p)t_{\mathrm{tr}}.
 \label{eq:transient-time}
\end{equation}
در رابطه \eqref{eq:transient-time}، $t_{\mathrm{prot,eff}}$ زمان مؤثر مورد انتظار وقفه بار محافظت‌شده، $t_{\mathrm{reclose}}$ زمان اتصال مجدد، $p$ احتمال موفقیت وصل و $t_{\mathrm{tr}}$ زمان وقفه گذرا است. این رابطه ساده‌سازی مدل امید ریاضی زمان وقفه است.

در تنظیم نمونه مستند پروژه، مقادیر آموزشی $t_{\mathrm{tr}}=1$ دقیقه، $t_{\mathrm{reclose}}=0.25$ دقیقه و $p=0.70$ استفاده شده‌اند؛ بنابراین:
\begin{equation}
 t_{\mathrm{prot,eff}}=0.25+(1-0.70)\times1=0.55\unit{min}.
 \label{eq:transient-numeric}
\end{equation}
در معادله \eqref{eq:transient-numeric}، مدت وقفه موثر ۰٫۵۵ دقیقه از تنظیم‌های نمونه به دست آمده است؛ این اعداد باید در مطالعه واقعی با تنظیمات ریکلوزر و داده رخداد جایگزین شوند.

\section{\lr{EENS} گذرا}
اگر بخشی از بار زیر پوشش ریکلوزر و بخشی خارج از آن باشد، انرژی مورد انتظار گذرا برای زون $i$ در مدل نمونه به‌شکل زیر است:
\begin{equation}
 EENS_{\mathrm{tr},i}=\frac{\lambda_{\mathrm{tr},i}}{60}\left[(P_{\mathrm{tot}}-P_{\mathrm{prot},i})t_{\mathrm{tr}}+P_{\mathrm{prot},i}\left(t_{\mathrm{reclose}}+(1-p)t_{\mathrm{tr}}\right)\right].
 \label{eq:transient-eens}
\end{equation}
در رابطه \eqref{eq:transient-eens}، $\lambda_{\mathrm{tr},i}$ نرخ سالانه خطاهای گذرای زون $i$، $P_{\mathrm{tot}}$ کل بار متأثر، $P_{\mathrm{prot},i}$ بار تحت پوشش ریکلوزر، $t_{\mathrm{tr}}$ زمان وقفه گذرا، $t_{\mathrm{reclose}}$ زمان وصل مجدد و $p$ احتمال موفقیت است. تقسیم بر ۶۰ تبدیل \lr{kW-minute} به \lr{kWh/year} می‌کند.

\subsection{مثال حساسیت با داده \file{5.xlsx}}
در \file{5.xlsx}، رخداد \lr{3429483} نمونه مستند خطای گذرا و وصل ریکلوزر است. اگر در یک تحلیل جداگانه تعداد ۸۹ رخداد با علت خام «عیب گذرا» در همین فایل به‌عنوان زیرمجموعه گذرا پذیرفته شود، نرخ ساده آن بر اساس مواجهه ۴٫۹۲۰۵۴۸ سال برابر $18.087$ رخداد بر سال خواهد بود. این نرخ، نرخ خام طبقه‌بندی‌شده بر اساس متن علت است و پیش از کاربرد نهایی باید با نمونه‌خوانی و قواعد رخداد تطبیق داده شود.

\assumption برای نمایش عددی، $P_{\mathrm{tot}}=3073.5$ \lr{kW} را از سناریوی بار بخش ۵.۲ می‌گیریم و فرض می‌کنیم نیمی از بار تحت پوشش ریکلوزر است؛ یعنی $P_{\mathrm{prot}}=1536.75$ \lr{kW}. با تنظیم‌های نمونه بخش ۸.۲:
\begin{equation}
 EENS_{\mathrm{tr}}=\frac{18.087}{60}\left[1536.75(1)+1536.75(0.55)\right]
 \simeq718.1\unit{kWh/year}.
 \label{eq:f5-transient-example}
\end{equation}
در معادله \eqref{eq:f5-transient-example}، نرخ گذرا ۱۸٫۰۸۷ رخداد بر سال، بار کل ۳۰۷۳٫۵ \lr{kW}، بار تحت حفاظت ۱۵۳۶٫۷۵ \lr{kW} و زمان‌های وقفه ۱ و ۰٫۵۵ دقیقه هستند. این مثال از نظر آموزشی نشان می‌دهد که اثر ریکلوزر تابع سهم بار تحت حفاظت است؛ نتیجه واقعی وابسته به تنظیمات و محل نصب است.

\section{زمان مرده و شات‌های ریکلوزر}
زمان مرده فاصله زمانی بین بازشدن ریکلوزر و صدور فرمان وصل مجدد است. در برنامه چندشاتی، ممکن است زمان مرده برای شات اول و شات‌های بعدی متفاوت باشد. اگر ریکلوزر دو شاته باشد، باید مشخص شود که کدام رخدادها با شات اول رفع شده‌اند و کدام رخدادها تا شات بعدی یا قفل نهایی پیش رفته‌اند. ثبت «وصل موفق» در متن حادثه به‌تنهایی برای استخراج توالی دقیق شات‌ها کافی نیست.

\warn مدت‌های ۲ و ۳۰ ثانیه‌ای که ممکن است در تنظیمات عملیاتی ریکلوزر مشاهده شوند، نباید بدون مستندات تنظیمی به کل فیدر یا تمام رخدادها تعمیم داده شوند. برای محاسبه معتبر، منحنی و توالی واقعی ریکلوزر، اطلاعات رخدادنگار و وضعیت کنتاکت‌ها نیاز است.

\summary{۸}
خطای گذرا و دائمی رفتار یکسان ندارند. در مدل ریکلوزر، احتمال وصل موفق و زمان وقفه مؤثر تعیین‌کننده‌اند. مثال \file{5.xlsx} نشان می‌دهد داده علت و شرح می‌تواند برای ساخت نرخ گذرا استفاده شود، ولی این نرخ باید از نظر مهندسی بازبینی شود.

\exercise
اگر $p=0.8$، $t_{\mathrm{reclose}}=0.2$ دقیقه و $t_{\mathrm{tr}}=1$ دقیقه باشد، زمان وقفه مؤثر بار محافظت‌شده چقدر است؟
\answer $0.2+(1-0.8)\times1=0.4$ دقیقه.

\chapter{فرمول‌بندی ریاضی مسئله جایابی}
\section{متغیرهای تصمیم و تابع هدف}
فرض کنیم مجموعه نقاط کاندید برای نصب تجهیزات $\mathcal{C}$ باشد. متغیر دودویی تعیین می‌کند در نقطه $j$ تجهیز نصب شود یا خیر:
\begin{equation}
 x_j=\begin{cases}1,&\text{اگر در نقطه }j\text{ تجهیز نصب شود},\\0,&\text{در غیر این صورت.}\end{cases}
 \label{eq:binary}
\end{equation}
در معادله \eqref{eq:binary}، $x_j$ متغیر دودویی تصمیم و $j$ شاخص نقطه کاندید است.

تابع هدف اصلی کمینه‌کردن \lr{EENS} سالانه است:
\begin{equation}
 \min_{\bm{x}}\;EENS(\bm{x})
 =\min_{\bm{x}}\sum_{i=1}^{n}EENS_i(\bm{x}).
 \label{eq:objective}
\end{equation}
در رابطه \eqref{eq:objective}، $\bm{x}$ بردار تمام متغیرهای تصمیم، $n$ تعداد بخش‌ها و $EENS_i(\bm{x})$ انرژی مورد انتظار بخش $i$ در آرایش انتخاب‌شده تجهیزات است.

\section{محدودیت تعداد و بودجه}
اگر حداکثر تعداد تجهیزات قابل نصب $K$ باشد:
\begin{equation}
 \sum_{j\in\mathcal{C}}x_j\le K.
 \label{eq:budget}
\end{equation}
در رابطه \eqref{eq:budget}، $\mathcal{C}$ مجموعه نقاط کاندید و $K$ سقف تعداد تجهیزات مجاز است. می‌توان به‌جای تعداد، محدودیت بودجه مالی و هزینه هر تجهیز را لحاظ کرد.

\begin{equation}
 \sum_{j\in\mathcal{C}}c_jx_j\le B.
 \label{eq:cost}
\end{equation}
در رابطه \eqref{eq:cost}، $c_j$ هزینه نصب تجهیز در نقطه $j$ و $B$ بودجه کل است؛ هر دو باید با واحد پولی یکسان بیان شوند.

\section{قیود مهندسی کاندیدا}
هر نقطه کاندید باید از نظر توپولوژی، سطح ولتاژ، قابلیت نصب، دسترسی، ایمنی و هماهنگی حفاظتی معتبر باشد. در منطق پیاده‌سازی بررسی‌شده، نقاطی مانند منبع، نقاط دارای تجهیزات موجود یا نقاط کم‌عمق ممکن است از مجموعه کاندیدا حذف شوند؛ حداقل عمق، حداقل بار زیرشاخه و سقف نسبت بار زیرشاخه به کل فیدر نیز می‌توانند فیلتر شوند. مقادیر پیش‌فرض نرم‌افزار، تنظیمات مدل‌اند و جایگزین ضوابط بهره‌برداری شرکت توزیع نیستند.
در تنظیم پیش‌فرض پیاده‌سازی بررسی‌شده، حداقل عمق کاندیدا ۳، بار زیرشاخه بزرگ‌تر از صفر و نسبت بار زیرشاخه به کل فیدر حداکثر ۰٫۸ در نظر گرفته می‌شود. این فیلترها برای حذف نقاط کم‌اثر یا نامناسب در فضای جست‌وجو هستند؛ مقدار نهایی آن‌ها باید متناسب با سیاست شرکت، ساختار فیدر و محدودیت‌های فنی تنظیم شود.

\section{روش حریصانه و معنای آن}
حل دقیق مسئله برای شبکه‌های بزرگ ممکن است پرهزینه باشد. روش حریصانه در هر گام، تجهیزی را انتخاب می‌کند که بیشترین کاهش \lr{EENS} افزایشی را ایجاد کند. اگر مجموعه فعلی تجهیزات $S$ باشد، منفعت افزودن گزینه $j$ چنین تعریف می‌شود:
\begin{equation}
 \Delta_j(S)=EENS(S)-EENS(S\cup\{j\}).
 \label{eq:marginal}
\end{equation}
در رابطه \eqref{eq:marginal}، $EENS(S)$ انرژی مورد انتظار با مجموعه تجهیزات فعلی، $EENS(S\cup\{j\})$ انرژی پس از افزودن تجهیز در نقطه $j$ و $\Delta_j(S)$ کاهش حاشیه‌ای EENS است.

الگوریتم در هر تکرار گزینه‌ای با بیشترین $\Delta_j(S)$ مثبت را انتخاب می‌کند و تا رسیدن به محدودیت تعداد/بودجه یا نبود بهبود ادامه می‌دهد. این راهبرد سریع و قابل فهم است، اما تضمین نمی‌کند جواب بهینه سراسری را پیدا کند؛ زیرا منفعت دو تجهیز ممکن است به یکدیگر وابسته باشد.

\section{مقایسه طرح‌ها و معیار تصمیم}
طرحی با کاهش \lr{EENS} بیشتر لزوماً اقتصادی‌تر نیست. اگر هزینه‌ها متفاوت‌اند، می‌توان نسبت کاهش \lr{EENS} به هزینه را نیز بررسی کرد:
\begin{equation}
 \eta_j=\frac{\Delta_j(S)}{c_j}.
 \label{eq:benefit-cost}
\end{equation}
در رابطه \eqref{eq:benefit-cost}، $\eta_j$ شاخص کاهش \lr{EENS} به ازای واحد هزینه، $\Delta_j(S)$ کاهش حاشیه‌ای و $c_j$ هزینه تجهیز است. این نسبت برای غربالگری مفید است، اما جایگزین تحلیل اقتصادی چرخه عمر و محدودیت‌های فنی نیست.

\summary{۹}
جایابی یک مسئله بهینه‌سازی گسسته است. متغیرهای دودویی محل نصب را نشان می‌دهند و تابع هدف \lr{EENS} را کمینه می‌کند. روش حریصانه با انتخاب بیشترین کاهش حاشیه‌ای کار می‌کند، ولی ممکن است به بهینه سراسری نرسد.

\exercise
اگر طرح A، \lr{EENS} را از ۱۰۰۰ به ۸۰۰ \lr{kWh/year} و طرح B آن را از ۱۰۰۰ به ۷۰۰ \lr{kWh/year} کاهش دهد، کدام کاهش \lr{EENS} بیشتری دارد؟
\answer طرح B؛ کاهش آن ۳۰۰ \lr{kWh/year} و کاهش طرح A برابر ۲۰۰ \lr{kWh/year} است. برای تصمیم اقتصادی هنوز هزینه‌ها و امکان‌پذیری فنی لازم‌اند.

\chapter{مطالعه موردی یکپارچه فیدر F5}
\section{گردآوری ورودی‌های واقعی و فرضی}
ورودی‌های واقعی از \file{5.xlsx} و لایه \lr{GIS} عبارت‌اند از: ۳۴۲ رکورد خام، ۲۱۵ رخداد تصادفی، ۱۱۳ رخداد قابل کاهش، مواجهه ۱۷۹۶ روز، $\lambda=43.6943$ رخداد بر سال، $r=19.051$ دقیقه، ۷۸ ترانس در \code{ELWDTR}، مجموع ظرفیت نامی ۶۸۳۰ \lr{kVA} و ۲۰۴۳ قطعه خط در \code{ELWOML} با جمع طول \code{Shape\_STLe} برابر ۱۲۳٫۱۵۹ \lr{km}.

فرض‌های صرفاً آموزشی در این مثال عبارت‌اند از $LF=0.50$، $PF=0.90$، $T_{\mathrm{rep}}=50$ دقیقه، $T_{\mathrm{restore}}=20$ دقیقه، تقسیم برابر بار و طول به دو زون و ظرفیت پشتیبان نامحدود. هر مقدار از این گروه باید در مطالعه اجرایی با داده واقعی جایگزین شود.

\section{گام اول: نرخ رخداد}
\begin{equation}
 \lambda_{F5}=\frac{215}{1796/365}=43.6943\;\mathrm{event/year}.
 \label{eq:case-step1}
\end{equation}
در معادله \eqref{eq:case-step1}، ۲۱۵ تعداد رخدادهای تصادفی و ۱۷۹۶ روز مدت مواجهه است.

\section{گام دوم: نرخ قابل کاهش}
\begin{equation}
 \lambda_{\mathrm{red},F5}=\frac{113}{1796/365}=22.9649\;\mathrm{event/year}.
 \label{eq:case-step2}
\end{equation}
در معادله \eqref{eq:case-step2}، ۱۱۳ تعداد رخدادهای طبقه «قابل کاهش با کلیدگذاری» است. نسبت این نرخ به نرخ تصادفی نشان می‌دهد همه رخدادهای تصادفی با کلیدگذاری قابل کاهش نیستند.

\section{گام سوم: بار برآوردی}
\begin{equation}
 P_{F5,\mathrm{est}}=6830\times0.50\times0.90=3073.5\unit{kW}.
 \label{eq:case-step3}
\end{equation}
در معادله \eqref{eq:case-step3}، ظرفیت نامی ترانس واقعی است، اما ضریب بار و ضریب توان فرض آموزشی‌اند؛ بنابراین $P_{F5,\mathrm{est}}$ تخمین سناریویی است.

\section{گام چهارم: عدم‌دسترس‌پذیری}
\begin{equation}
 U_{F5}=\frac{43.6943\times19.051}{60}=13.8737\unit{h/year}.
 \label{eq:case-step4}
\end{equation}
در معادله \eqref{eq:case-step4}، $\lambda$ نرخ رخداد تصادفی و $r$ متوسط مدت خاموشی تصادفی است.

\section{گام پنجم: مقایسه \lr{EENS} پیش و پس از تقسیم فرضی}
مقادیر سناریوی بدون جداسازی و دو زون در معادلات \eqref{eq:f5-base-eens} و \eqref{eq:f5-new-eens} محاسبه شدند. کاهش درصدی:
\begin{equation}
 \Delta EENS=100\frac{111912.1-78338.5}{111912.1}=30.0\%.
 \label{eq:case-step5}
\end{equation}
در رابطه \eqref{eq:case-step5}، صورت کسر کاهش \lr{EENS} و مخرج مقدار مبناست. به علت فرض مساوی بودن بار، طول و زمان بازیابی، این درصد یک نتیجه تحلیلی برای سناریوی ساده است؛ نباید آن را به‌عنوان درصد بهبود واقعی فیدر کوثر گزارش کرد.

\section{چه داده‌هایی برای تبدیل این مثال به طرح اجرایی لازم است؟}
برای طراحی واقعی حداقل موارد زیر لازم‌اند: (۱) بار اکتیو هر ترانس یا هر زون در بازه معتبر، (۲) توپولوژی الکتریکی متصل و وضعیت عادی کلیدها، (۳) محل و نوع دقیق هر تجهیز حفاظتی، (۴) زمان تعمیر و زمان مانور واقعی، (۵) ظرفیت انتقال مسیرهای جایگزین و محدودیت ولتاژ/جریان، (۶) تفکیک مطمئن رخداد گذرا و دائمی، و (۷) واحد تأییدشده انرژی ثبت‌شده. بدون این داده‌ها، محاسبات \lr{EENS} باید سناریویی و نه قطعی تلقی شوند.

\summary{۱۰}
مثال یکپارچه \code{F5} نشان داد چگونه از فایل خاموشی به نرخ رخداد و سپس به تخمین \lr{EENS} می‌رسیم. داده‌های پایه نرخ و ظرفیت ترانس واقعی‌اند، اما بار اکتیو و تقسیم زون در سناریوی آموزشی فرض شده‌اند. شفافیت این تمایز، شرط استفاده علمی از خروجی است.

\exercise
اگر در مثال دو زون، زمان بازیابی بار سالم به جای ۲۰ دقیقه به ۱۰ دقیقه کاهش یابد و سایر فرض‌ها ثابت بمانند، \lr{EENS} جدید در سناریوی دو زون چقدر می‌شود؟
\answer از رابطه \eqref{eq:f5-new-eens}، به جای ۲۰ مقدار ۱۰ قرار می‌دهیم:
$2\times(43.6943/2)/60\times[50(1536.75)+10(1536.75)]\simeq67{,}170.4$ \lr{kWh/year}.

\chapter{تحلیل حساسیت، اعتبارسنجی و محدودیت‌ها}
\section{حساسیت به پوشش آماری}
اگر مخرج مواجهه کمتر از واقع برآورد شود، نرخ رخداد بزرگ‌تر به دست می‌آید؛ اگر ماه‌های بدون داده به اشتباه ماه‌های بدون رخداد تلقی شوند، نرخ کمتر از واقع خواهد شد. بنابراین نمودار ماه‌های پوشش‌دار، تاریخ آخرین ثبت، وقفه سامانه ثبت و تغییرات محدوده فیدر باید کنترل شوند.

\section{تفکیک نرخ قابل کاهش از نرخ کل در مدل EENS}
در فایل‌های خاموشی، $\lambda$ شامل سه گروه «قابل کاهش با کلیدگذاری»، «نیازمند تعمیر یا تعویض تجهیز» و «ناشی از هماهنگی حفاظتی» است؛ اما فقط $\lambda_{\mathrm{red}}$ به‌طور مستقیم نماینده رخدادهای کلاس قابل کاهش با کلیدگذاری است. بنابراین باید مشخص شود در محاسبه \lr{EENS} چه فرضی اعمال می‌شود. اگر تمام نرخ $\lambda$ به مدل جایابی داده شود، مدل ساده فرض می‌کند اثر جداسازی روی بار سالم برای تمام رخدادهای تصادفی قابل اعمال است. اگر فقط $\lambda_{\mathrm{red}}$ استفاده شود، اثر قابل انتساب به رخدادهای کلاس اول محدودتر و محافظه‌کارانه‌تر خواهد بود.

\begin{equation}
 \lambda_{\mathrm{other}}=\lambda-\lambda_{\mathrm{red}},\qquad
 EENS_{\mathrm{all}}(S)=EENS_{\mathrm{red}}(S;\lambda_{\mathrm{red}})+EENS_{\mathrm{other}}(S).
 \label{eq:lambda-split}
\end{equation}
در رابطه \eqref{eq:lambda-split}، $\lambda_{\mathrm{other}}$ نرخ رخدادهای تصادفیِ خارج از طبقه قابل کاهش، $S$ آرایش تجهیزات انتخاب‌شده، $EENS_{\mathrm{red}}$ انرژی مورد انتظار ناشی از رخدادهای قابل کاهش و $EENS_{\mathrm{other}}$ انرژی رخدادهایی است که اثر کلیدگذاری بر آن‌ها باید جداگانه مدل شود. برای فیدر \code{F5}، مقدار $\lambda_{\mathrm{other}}=43.6943-22.9649=20.7294$ رخداد بر سال است.

در کد محاسبه \lr{EENS} بررسی‌شده، تابع محاسبه از یک ورودی نرخ خرابی کلی استفاده می‌کند و به‌تنهایی طبقه‌بندی رخداد را به چند نرخ مستقل تفکیک نمی‌کند. بنابراین انتخاب نرخ ورودی، فرض کلیدی مدل است. برای مطالعه دقیق‌تر، نرخ رخدادهای قابل کاهش و نرخ سایر رخدادها باید جداگانه وارد شوند تا کاهش \lr{EENS} به‌اشتباه به تمام گروه‌ها نسبت داده نشود.

\section{حساسیت به صفرهای مدت خاموشی}
در \file{5.xlsx}، متوسط مدت رخدادهای تصادفی با حفظ رخدادهای صفر برابر ۱۹٫۰۵۱ دقیقه است؛ میانگین حساسیت با حذف صفرها ۳۳٫۵۷۴ دقیقه گزارش شده است. اختلاف قابل توجه نشان می‌دهد رخدادهای صفر دقیقه‌ای در مجموعه اثرگذارند. کدام میانگین برای مطالعه مناسب است به تعریف شاخص وابسته است: برای مدت متوسط هر رخداد واقعی، صفرهای معتبر باید باقی بمانند؛ برای زمان بازیابی رخدادهای دارای وقفه قابل اندازه‌گیری، میانگین بدون صفر می‌تواند شاخص مکمل باشد. این دو باید جداگانه گزارش شوند.

\section{حساسیت به بار و ضریب توان}
بار برآوردی به شکل مستقیم با $LF$ و $PF$ متناسب است. اگر بار واقعی ۲۰ درصد بیشتر از فرض باشد، در مدل خطی و با سایر شرایط ثابت، \lr{EENS} نیز ۲۰ درصد بیشتر می‌شود. به همین دلیل، استفاده از ظرفیت نامی ترانس به‌عنوان بار واقعی می‌تواند خطای بزرگی ایجاد کند.

\section{حساسیت به فرض نرخ طولی}
مدل طول‌محور معادله \eqref{eq:length-lambda} فرض می‌کند نرخ خرابی به ازای کیلومتر در تمام خط یکسان است. این فرض ممکن است برای مسیرهای با پوشش گیاهی متفاوت، هادی‌های فرسوده، مناطق بادخیز، آلودگی، صاعقه یا شبکه‌های کابل زمینی برقرار نباشد. در صورت دسترسی به داده مکانی دقیق‌تر، می‌توان ضریب نرخ را بر اساس کلاس هادی، سن، محیط یا سابقه خرابی اصلاح کرد.

\section{محدودیت اطلاعات و ریسک تفسیر}
\begin{itemize}
\item فایل‌های خاموشی ستون تعداد مشترکین قطع‌شده یا توان قطع‌شده هر رخداد را ندارند؛ بنابراین محاسبه مستقیم \lr{SAIFI}، \lr{SAIDI} و \lr{EENS} دقیق از همین فایل‌ها ممکن نیست.
\item واحد ستون انرژي توزيع نشده در فایل منبع تصریح نشده است؛ تا زمان تأیید، آن را وارد جمع \lr{kWh} نکنید.
\item ستون مستقلی برای «برنامه‌ریزی‌شده/اضطراری» وجود ندارد؛ شناسایی خاموشی برنامه‌ریزی‌شده به قواعد متنی وابسته است.
\item نام تجهیز عمل‌کرده همیشه محل دقیق آن را مشخص نمی‌کند؛ به‌ویژه «کات اوت سرخط» به بازبینی کارشناس نیاز دارد.
\item داده \lr{GIS} ظرفیت نامی ترانس و هندسه خطوط را می‌دهد، نه توان مصرفی واقعی در ساعت حادثه.
\item نتیجه بهینه‌سازی حریصانه، جواب پیشنهادی الگوریتم است و باید از نظر بهره‌برداری، حفاظت، ایمنی و هزینه اعتبارسنجی شود.
\end{itemize}

\section{فهرست کنترل پیش از انتشار نتیجه مهندسی}
\begin{enumerate}
\item تعداد رکوردها و شماره حادثه‌های تکراری کنترل شود.
\item تبدیل تاریخ شمسی و زمان \lr{Excel} بازبینی شود.
\item ماه‌های دارای داده و ماه‌های مفقود مستند شوند.
\item علت، تجهیز و شرح رخدادهای اثرگذار توسط متخصص حفاظت نمونه‌خوانی شوند.
\item رخدادهای داده ناکافی در نرخ نهایی وارد نشوند تا تعیین تکلیف شوند.
\item اتصال \lr{GIS} به فیدر و هماهنگی شناسه‌ها کنترل شود.
\item بار هر زون با اندازه‌گیری یا برآورد مستند تعیین شود.
\item سناریوهای زمان تعمیر، زمان بازیابی و ظرفیت پشتیبان بررسی شوند.
\item \lr{EENS} مبنا و طرح جدید با مجموعه ورودی یکسان محاسبه شوند.
\item محل نصب نهایی از نظر هماهنگی حفاظتی، ایمنی و امکان مانور تأیید شود.
\end{enumerate}

\summary{۱۱}
تحلیل قابلیت اطمینان نسبت به کیفیت داده و فرض‌های مدل حساس است. نرخ‌ها و میانگین‌ها باید همراه با پوشش، مقدارهای صفر، رخدادهای ناکافی و فرض بار گزارش شوند. عدد \lr{EENS} بدون واحد معتبر بار، مدت و انرژی قابل اتکا نیست.

\exercise
چرا نمی‌توان از ستون انرژي توزيع نشده مستقیماً مجموع \lr{kWh} گرفت؟
\answer چون واحد این فیلد در فایل خام مشخص نشده است. ابتدا باید از مالک سامانه یا فرهنگ داده، واحد و روش محاسبه آن تأیید شود.

\chapter{واژه‌نامه نمادها و شاخص‌ها}
\begin{longtable}{p{2.2cm}p{4.2cm}p{8.0cm}}
\caption{تعریف نمادهای پرکاربرد در روابط این راهنما}\label{tab:symbols}\\
\toprule نماد & واحد & تعریف \\
\midrule
\endfirsthead
\toprule نماد & واحد & تعریف \\
\midrule
\endhead
$G$ & - & گراف مدل شبکه.\\
$V$ & - & مجموعه رأس‌های گراف.\\
$E$ & - & مجموعه یال‌ها/قطعات خط.\\
$N_{\mathrm{random}}$ & رخداد & تعداد رخدادهای تصادفی واردشده به مدل.\\
$N_{\mathrm{reducible}}$ & رخداد & تعداد رخدادهای قابل کاهش با کلیدگذاری.\\
$Y_{\mathrm{exp}}$ & سال & مدت مواجهه معتبر داده.\\
$T_{\mathrm{exp}}$ & روز & مجموع روزهای مواجهه معتبر.\\
$\lambda$ & رخداد/سال & نرخ سالانه رخدادهای تصادفی.\\
$\lambda_{\mathrm{red}}$ & رخداد/سال & نرخ سالانه رخدادهای قابل کاهش با کلیدگذاری.\\
$\lambda_i$ & رخداد/سال & نرخ تخصیص‌یافته به زون یا بخش $i$.\\
$r$ & دقیقه/رخداد & متوسط مدت خاموشی رخدادهای منتخب.\\
$t_j$ & دقیقه & مدت رخداد $j$.\\
$U$ & ساعت/سال & عدم‌دسترس‌پذیری/مدت خاموشی مورد انتظار سالانه در مدل ساده.\\
$ENS$ & \lr{kWh} & انرژی تأمین‌نشده در رخداد یا بازه مشخص.\\
$EENS$ & \lr{kWh/year} & انرژی تأمین‌نشده مورد انتظار سالانه.\\
$P_i$ & \lr{kW} & بار اکتیو زون یا بخش $i$.\\
$P_{\mathrm{lost}}$ & \lr{kW} & توان قطع‌شده در زمان مورد بررسی.\\
$P_{\mathrm{healthy}}$ & \lr{kW} & بار سالمی که پس از خطا می‌تواند بالقوه از مسیر دیگر تغذیه شود.\\
$P_{\mathrm{restored}}$ & \lr{kW} & بار سالمی که واقعاً پس از مانور بازیابی می‌شود.\\
$P_{\mathrm{remaining}}$ & \lr{kW} & بار سالمی که تا زمان تعمیر بی‌برق می‌ماند.\\
$S_j$ & \lr{kVA} & ظرفیت نامی ترانسفورماتور $j$.\\
$LF$ & - & ضریب بار.\\
$PF$ & - & ضریب توان.\\
$L_i$ & \lr{km} یا \lr{m} & طول بخش $i$ با واحد یکسان در کل محاسبه.\\
$T_{\mathrm{rep}}$ & دقیقه & زمان تعمیر تجهیز معیوب.\\
$T_{\mathrm{restore}}$ & دقیقه & زمان مانور و بازیابی بار سالم.\\
$T_{\mathrm{rest}}$ & دقیقه & زمان مؤثر بازیابی، محدودشده به زمان تعمیر.\\
$C_{\mathrm{backup}}$ & \lr{kW} & ظرفیت بار قابل انتقال در مسیر پشتیبان.\\
$x_j$ & - & متغیر دودویی نصب تجهیز در نقطه کاندید $j$.\\
$K$ & تجهیز & حداکثر تعداد تجهیزات قابل نصب.\\
$c_j$ & واحد پول & هزینه تجهیز و نصب در نقطه $j$.\\
$B$ & واحد پول & بودجه مجاز.\\
$\Delta_j(S)$ & \lr{kWh/year} & کاهش حاشیه‌ای EENS با افزودن گزینه $j$ به مجموعه $S$.\\
$p$ & - & احتمال موفقیت وصل مجدد ریکلوزر.\\
$t_{\mathrm{tr}}$ & دقیقه & مدت وقفه گذرای مدل‌شده.\\
$t_{\mathrm{reclose}}$ & دقیقه & زمان وصل مجدد/بازیابی بار تحت حفاظت.\\
$\lambda_{\mathrm{tr},i}$ & رخداد/سال & نرخ خطاهای گذرای زون $i$.\\
\bottomrule
\end{longtable}

\chapter{منابع داده و ملاحظات بازتولیدپذیری}
\begin{enumerate}
\item فایل خاموشی فیدر ۳: \file{3.xlsx}، شیت \code{Outages}، آرشیو \file{kosar.rar}.
\item فایل خاموشی فیدر ۴: \file{4.xlsx}، شیت \code{Outages}، آرشیو \file{kosar.rar}.
\item فایل خاموشی فیدر ۵: \file{5.xlsx}، شیت \code{Outages}، آرشیو \file{kosar.rar}.
\item فایل خاموشی فیدر ۶: \file{6.xlsx}، شیت \code{Outages}، آرشیو \file{kosar.rar}.
\item فایل خاموشی فیدر ۷: \file{7.xlsx}، شیت \code{Outages}، آرشیو \file{kosar.rar}.
\item فایل تک‌ردیفی: \file{بدون نام فیدر.xlsx}.
\item فایل مکانی داخل آرشیو: \file{Kosar.zip}، مسیر \file{Export\_\allowbreak 1404\_\allowbreak 08\_\allowbreak 17\_\allowbreak 12\_\allowbreak 26\_\allowbreak 05}، از جمله لایه‌های \code{ELWDTR}، \code{ELWOML}، \code{ELWUML}، \code{ELWCBR}، \code{ELWDCS}، \code{ELWFCT} و \code{ELWFWT}.
\item بسته نرم‌افزار: \file{SwitchLoc-\allowbreak arena-\allowbreak fad3d399-\allowbreak v1.22.0.zip}.
\item مستندات پردازش خاموشی: \file{OUTAGE\_PIPELINE\_FA.md}.
\item مستندات جایابی تجهیزات: \file{SWITCH\_PLACEMENT\_FA.md}.
\item مستندات مقایسه \lr{EENS}: \file{COMPARISON\_EENS\_FA.md}.
\item مسیر کدهای قابلیت اطمینان: \file{src/\allowbreak switchloc/\allowbreak reliability/}.
\end{enumerate}

\section*{یادداشت نهایی}
اعداد این راهنما برای آموزش و بازتولید مسیر محاسباتی از فایل‌های پیوست‌شده استخراج شده‌اند. پیش از استفاده در گزارش رسمی شرکت توزیع یا تصمیم سرمایه‌گذاری، باید نسخه داده، واحد فیلدها، تعریف رخداد تصادفی، بار واقعی، پارامترهای حفاظت و محدودیت‌های بهره‌برداری با متولیان مربوط تأیید شوند. به‌ویژه سناریوی EENS فصل ۵ و مطالعه موردی فصل ۱۰ عمداً فرض‌های آموزشی آشکار دارند و نتیجه محاسبه رسمی شبکه محسوب نمی‌شوند.

\backmatter
\end{document}
